\section{Appendix: Complete Proofs of Theorems}
\label{sec:appendix}

\subsection{Proof of Theorem 1: Lipschitz Stability of the Input Gradient}
Here we provide the detailed algebraic derivation of the input gradient perturbation bound stated in Theorem~\ref{thm:lipschitz}.

\begin{proof}
Let the class logit score $S_c(\bx)$ for an input vector $\bx \in \mathbb{R}^d$ under an $L$-layer neural network parameterised by weights $\bm{\theta} = \{\bW_1, \bW_2, \dots, \bW_L\}$ be written as:
\begin{equation}
S_c(\bx) = \bW_L^{\top} \sigma\Bigl(\bW_{L-1} \sigma\bigl(\dots \sigma(\bW_1 \bx)\bigr)\Bigr) ,
\end{equation}
where $\sigma$ is a element-wise activation function. The input gradient attribution map $\bg = \nabla_{\bx} S_c(\bx)$ is obtained via the chain rule as a product of layer weight matrices and activation derivative Jacobians:
\begin{equation}
\bg = \bW_1^{\top} D_1 \bW_2^{\top} D_2 \dots \bW_{L-1}^{\top} D_{L-1} \bW_L ,
\label{eq:gradchain}
\end{equation}
where $D_l = \mathrm{diag}\bigl(\sigma'(\bm{z}_l)\bigr) \in \mathbb{R}^{n_l \times n_l}$ is the diagonal Jacobian of the activation function at layer $l$, evaluated at the pre-activation vector $\bm{z}_l$. Since $\sigma$ is typically a standard activation function (such as ReLU, LeakyReLU, or Sigmoid), its derivative is bounded, satisfying:
\begin{equation}
\norm{D_l}_2 = \max_i |\sigma'([\bm{z}_l]_i)| \le 1 \quad \forall l \in \{1, \dots, L-1\} .
\end{equation}

Now consider a perturbation applied to each layer's weights, parameterised as $\tilde{\bW}_l = \bW_l + \bE_l$. The perturbed input gradient is:
\begin{equation}
\tilde{\bg} = (\bW_1 + \bE_1)^{\top} \tilde{D}_1 (\bW_2 + \bE_2)^{\top} \tilde{D}_2 \dots (\bW_{L-1} + \bE_{L-1})^{\top} \tilde{D}_{L-1} (\bW_L + \bE_L) ,
\label{eq:perturbedgrad}
\end{equation}
where $\tilde{D}_l = \mathrm{diag}\bigl(\sigma'(\tilde{\bm{z}}_l)\bigr)$ is the Jacobian evaluated under the perturbed pre-activations.

To bound the difference $\norm{\tilde{\bg} - \bg}_2$, we decompose the total perturbation into first-order and higher-order terms. First, under the local linear stability assumption (Assumption~\ref{ass:net}), the active activation regions are assumed to be fixed for small perturbations, meaning $\tilde{D}_l \approx D_l$. Expanding the product in \eqref{eq:perturbedgrad} and collecting terms linear in the perturbations $\bE_l$ gives:
\begin{equation}
\tilde{\bg} = \bg + \sum_{l=1}^L \bm{\delta}_l + \mathcal{R}(\{\bE_k\}) ,
\end{equation}
where each $\bm{\delta}_l$ represents the first-order perturbation term in which only the $l$-th weight matrix is replaced by its error $\bE_l$:
\begin{equation}
\bm{\delta}_l = \bW_1^{\top} D_1 \dots \bE_l^{\top} D_l \dots \bW_{L-1}^{\top} D_{L-1} \bW_L ,
\end{equation}
and $\mathcal{R}$ is the remainder containing all higher-order cross-terms $\bE_i \bE_j$ ($i \ne j$).

Applying the triangle inequality to the first-order approximation:
\begin{equation}
\norm{\tilde{\bg} - \bg}_2 \le \sum_{l=1}^L \norm{\bm{\delta}_l}_2 + \norm{\mathcal{R}}_2 .
\label{eq:triangle}
\end{equation}
Using the sub-multiplicativity of the spectral norm ($\norm{\bm{A}\bm{B}}_2 \le \norm{\bm{A}}_2 \norm{\bm{B}}_2$) and the bound $\norm{D_k}_2 \le 1$, the norm of each individual term $\bm{\delta}_l$ is bounded by:
\begin{align}
\norm{\bm{\delta}_l}_2 &= \norm{\bW_1^{\top} D_1 \dots \bE_l^{\top} D_l \dots \bW_{L-1}^{\top} D_{L-1} \bW_L}_2 \nonumber \\
&\le \left(\prod_{k \ne l} \norm{\bW_k}_2\right) \left(\prod_{k=1}^{L-1} \norm{D_k}_2\right) \norm{\bE_l}_2 \nonumber \\
&\le \left(\prod_{k \ne l} \norm{\bW_k}_2\right) \norm{\bE_l}_2 .
\end{align}
Substituting this back into \eqref{eq:triangle} yields the first-order bound:
\begin{equation}
\norm{\tilde{\bg} - \bg}_2 \le \sum_{l=1}^L \left(\prod_{k \ne l} \norm{\bW_k}_2\right) \norm{\bE_l}_2 + O(\norm{\Delta\btheta}_2^2) .
\end{equation}
By defining the gradient amplification factor $\Lambda_{\mathrm{S}} = \sum_{l=1}^L \prod_{k \ne l} \norm{\bW_k}_2$, and factoring out the maximum relative layer-wise error, we arrive at the desired result:
\begin{equation}
\norm{\tilde{\bg} - \bg}_2 \le \Lambda_{\mathrm{S}} \max_l \frac{\norm{\bE_l}_2}{\norm{\bW_l}_2} \norm{\bW_l}_2 + O(\norm{\Delta\btheta}_2^2) .
\label{eq:finalfirstorder}
\end{equation}

When the activation pattern varies ($\tilde{D}_l \ne D_l$), we bound the difference $\norm{\tilde{D}_l - D_l}_2$. Assuming the activation derivative $\sigma'$ is $\beta$-Lipschitz, we have:
\begin{equation}
\norm{\tilde{D}_l - D_l}_2 \le \beta \norm{\tilde{\bm{z}}_l - \bm{z}_l}_2 .
\end{equation}
Since the pre-activation shift $\norm{\tilde{\bm{z}}_l - \bm{z}_l}_2$ is bounded by $C \norm{\Delta\btheta}_2$ to first order, the perturbation term $\norm{\tilde{D}_l - D_l}_2$ contributes an additional error of order $O(\beta \norm{\Delta\btheta}_2^2)$, which is absorbed directly into the second-order curvature remainder. This completes the proof.
\end{proof}
